% appP_graph_e.tex -- generated by writing/tools/blueprint_gen.py 1.0 (level 3) on 2026-10-14 from sources 20021204b9; do not edit
\providecommand{\bpNum}[1]{\hyperref[#1]{\ref*{#1}}}%
\providecommand{\bpNumA}[1]{\hyperref[#1]{\textcolor{orange!65!black}{\ref*{#1}}}}%
\providecommand{\bpPage}[1]{\pageref*{#1}}%
\providecommand{\bpLab}[1]{{\ttfamily #1}}%
\providecommand{\bpBrk}{\allowbreak}%
\providecommand{\bpT}[2]{\facttagbox[#1]{\scriptsize #2}}%
\providecommand{\bpU}[1]{{\color{gray!70!black}#1}}%
\begin{figure}[tbp]
\centering
\begin{tikzpicture}[
  bp node/.style={draw,line width=0.45pt,rounded corners=1pt,font=\tiny\ttfamily,align=center,
                  inner xsep=1.5pt,inner ysep=1.1pt,minimum height=4mm},
  bp R/.style={bp node,fill=OliveGreen!15,draw=OliveGreen!60!black},
  bp A/.style={bp node,fill=RoyalBlue!15,draw=RoyalBlue!60!black},
  bp P/.style={bp node,fill=BurntOrange!15,draw=BurntOrange!70!black},
  bp C/.style={bp node,fill=Plum!15,draw=Plum!70!black},
  bp V/.style={bp node,fill=BrickRed!15,draw=BrickRed!70!black,double,double distance=0.6pt},
  bp cond/.style={dash pattern=on 1.6pt off 1pt},
  bp math/.style={bp node,fill=white,draw=black!60},
  bp def/.style={bp node,fill=white,draw=black!45,line width=0.3pt,rounded corners=0pt},
  bp conj/.style={bp node,ellipse,inner xsep=0pt,inner ysep=0pt,fill=Plum!15,draw=Plum!70!black},
  bp asm/.style={bp node,fill=orange!8,draw=orange!65!black,dash pattern=on 2pt off 1.2pt},
  bp stub/.style={bp node,font=\tiny,fill=gray!12,draw=gray!55,text=black!70},
  bp fold/.style={bp node,font=\tiny,fill=white,draw=gray!55,dash pattern=on 1pt off 1pt,text=black!70},
  bp unit/.style={font=\scriptsize,text=black!70,inner sep=1pt},
  bp sep/.style={draw=gray!40,line width=0.4pt,dash pattern=on 1pt off 1.5pt},
  bp key/.style={font=\scriptsize,anchor=west,inner sep=0pt},
  bp use/.style={draw=gray!60,line width=0.5pt,-{Stealth[length=2.6pt,width=2.2pt]}},
  bp stubuse/.style={draw=gray!45,line width=0.45pt,dash pattern=on 1.2pt off 1pt,-{Stealth[length=2.6pt,width=2.2pt]}},
  bp asmuse/.style={draw=orange!65!black,line width=0.5pt,dash pattern=on 2pt off 1.2pt,-{Stealth[length=2.6pt,width=2.2pt]}},
  bp neck/.style={draw=red!60!black,line width=1.1pt,-{Stealth[length=3.6pt,width=3.2pt]}},
  bp long/.style={opacity=0.5}]
  \node[bp unit] at (2.160,4.253) {第~14~章};
  \node[bp unit] at (5.270,4.253) {第~15~章};
  \node[bp unit] at (7.096,4.253) {第~16~章};
  \node[bp unit] at (9.152,4.253) {第~17~章};
  \node[bp unit] at (12.623,4.253) {第~18~章};
  \draw[bp sep] (0.663,4.113) -- (0.663,-3.983);
  \draw[bp sep] (3.888,4.113) -- (3.888,-3.983);
  \draw[bp sep] (6.422,4.113) -- (6.422,-3.983);
  \draw[bp sep] (7.770,4.113) -- (7.770,-3.983);
  \draw[bp sep] (10.997,4.113) -- (10.997,-3.983);
  \draw[bp sep] (14.274,4.113) -- (14.274,-3.983);
  \node[bp R] (n1) at (1.394,2.060) {\detokenize{tmd-sqrt-}\\\detokenize{identity}};
  \node[bp R,bp cond] (n2) at (2.926,1.583) {\detokenize{tmd-zero-}\\\detokenize{bin-soft}};
  \node[bp A] (n3) at (2.926,-0.200) {\detokenize{tmd-}\\\detokenize{factorization}};
  \node[bp A] (n4) at (2.926,0.692) {\detokenize{tmd-}\\\detokenize{evolution}};
  \node[bp A] (n5) at (2.926,2.475) {\detokenize{tmd-}\\\detokenize{small-b}};
  \node[bp R] (n6) at (2.926,3.489) {\detokenize{tmd-}\\\detokenize{sivers-}\\\detokenize{sign}};
  \node[bp A] (n7) at (2.926,-1.214) {\detokenize{tmd-}\\\detokenize{sivers-}\\\detokenize{observable}};
  \node[bp A] (n8) at (2.926,-2.352) {\detokenize{tmd-}\\\detokenize{collins-}\\\detokenize{universality}};
  \node[bp P] (n9) at (2.926,-3.489) {\detokenize{tmd-}\\\detokenize{twist3-}\\\detokenize{matching}};
  \node[bp R] (n10) at (4.677,0.480) {\detokenize{scet-}\\\detokenize{decoupling}};
  \node[bp A] (n11) at (5.864,0.953) {\detokenize{scet-}\\\detokenize{thrust}};
  \node[bp A] (n12) at (7.096,0.262) {\detokenize{jets-bms}};
  \node[bp P] (n13) at (7.096,2.555) {\detokenize{jets-sll}};
  \node[bp A] (n14) at (8.386,1.346) {\detokenize{smallx-}\\\detokenize{bfkl}};
  \node[bp P] (n15) at (9.919,1.671) {\detokenize{smallx-kt}};
  \node[bp A] (n16) at (9.919,1.021) {\detokenize{smallx-bk}};
  \node[bp conj] (n17) at (9.919,0.154) {\detokenize{smallx-}\\\detokenize{saturation}};
  \node[bp R,bp cond] (n18) at (11.844,1.346) {\detokenize{break-}\\\detokenize{measurement}};
  \node[bp V] (n19) at (13.402,1.746) {\detokenize{break-tmd}};
  \node[bp V] (n20) at (11.844,2.913) {\detokenize{break-}\\\detokenize{spacelike}};
  \node[bp A] (n21) at (11.844,0.575) {\detokenize{break-ddis}};
  \node[bp V] (n22) at (13.402,0.590) {\detokenize{break-}\\\detokenize{hh-}\\\detokenize{diffraction}};
  \node[bp A] (n23) at (11.844,-0.195) {\detokenize{break-}\\\detokenize{double-dy}};
  \node[bp fold] (n24) at (0.000,2.519) {+3};
  \node[bp stub] (n25) at (0.000,0.569) {第 3 章};
  \node[bp stub] (n26) at (0.000,1.869) {第 4 章};
  \node[bp stub] (n27) at (0.000,3.169) {第 5 章};
  \node[bp stub] (n28) at (0.000,1.219) {第 8 章};
  \node[bp stub] (n29) at (0.000,-0.081) {第 9 章};
  \node[bp stub] (n30) at (0.000,-0.731) {第 10 章};
  \node[bp stub] (n31) at (0.000,-1.381) {第 12 章};
  \node[bp fold] (n32) at (1.394,0.692) {+4};
  \node[bp fold] (n33) at (4.677,1.251) {+4};
  \node[bp fold] (n34) at (7.096,0.912) {+4};
  \node[bp fold] (n35) at (11.844,2.117) {+6};
  \node[bp stub] (n36) at (14.822,1.777) {附录 N};
  \node[bp fold] (n37) at (14.822,1.127) {+3};
  \node[bp stub] (n38) at (14.822,0.477) {附录 O};
  \node[bp fold] (n39) at (14.822,-0.173) {+3};
  \begin{scope}[on background layer]
  \draw[bp stubuse] (n24) -- (n3);
  \draw[bp stubuse] (n24) -- (n5);
  \draw[bp stubuse] (n24) -- (n6);
  \draw[bp stubuse,bp long] (n24) -- (n11);
  \draw[bp stubuse,bp long] (n24) -- (n21);
  \draw[bp stubuse,bp long] (n24) -- (n23);
  \draw[bp stubuse,bp long] (n30) -- (n21);
  \draw[bp stubuse] (n31) -- (n3);
  \draw[bp stubuse,bp long] (n31) -- (n23);
  \draw[bp stubuse] (n32) -- (n2);
  \draw[bp stubuse] (n32) -- (n3);
  \draw[bp use] (n1) -- (n2);
  \draw[bp use] (n1) -- (n4);
  \draw[bp use] (n1) -- (n5);
  \draw[bp use] (n1) -- (n6);
  \draw[bp use] (n2) to[bend left=35] (n3);
  \draw[bp use] (n3) to[bend left=35] (n7);
  \draw[bp use] (n3) to[bend left=35] (n8);
  \draw[bp use] (n3) to[bend left=35] (n9);
  \draw[bp use] (n6) to[bend left=35] (n7);
  \draw[bp stubuse] (n33) to[bend left=35] (n10);
  \draw[bp stubuse] (n33) -- (n11);
  \draw[bp use] (n10) -- (n11);
  \draw[bp stubuse] (n34) to[bend left=35] (n12);
  \draw[bp use] (n14) -- (n15);
  \draw[bp use] (n14) -- (n16);
  \draw[bp use] (n16) to[bend left=35] (n17);
  \draw[bp stubuse] (n35) -- (n19);
  \draw[bp use] (n18) -- (n19);
  \draw[bp use] (n21) -- (n22);
  \draw[bp stubuse] (n25) -- (n2);
  \draw[bp stubuse,bp long] (n25) -- (n10);
  \draw[bp stubuse,bp long] (n25) -- (n21);
  \draw[bp stubuse,bp long] (n25) -- (n23);
  \draw[bp stubuse,bp long] (n36) -- (n2);
  \draw[bp stubuse,bp long] (n37) -- (n2);
  \draw[bp stubuse,bp long] (n37) -- (n3);
  \draw[bp stubuse,bp long] (n37) -- (n4);
  \draw[bp stubuse,bp long] (n37) -- (n11);
  \draw[bp stubuse] (n37) -- (n21);
  \draw[bp stubuse] (n37) -- (n23);
  \draw[bp stubuse] (n26) -- (n2);
  \draw[bp stubuse] (n26) -- (n4);
  \draw[bp stubuse,bp long] (n26) -- (n11);
  \draw[bp stubuse,bp long] (n26) -- (n13);
  \draw[bp stubuse,bp long] (n26) -- (n18);
  \draw[bp stubuse,bp long] (n26) -- (n20);
  \draw[bp stubuse,bp long] (n26) -- (n21);
  \draw[bp stubuse,bp long] (n26) -- (n23);
  \draw[bp stubuse,bp long] (n38) -- (n3);
  \draw[bp stubuse,bp long] (n38) -- (n11);
  \draw[bp stubuse] (n38) -- (n18);
  \draw[bp stubuse] (n38) -- (n21);
  \draw[bp stubuse] (n38) -- (n23);
  \draw[bp stubuse,bp long] (n39) -- (n3);
  \draw[bp stubuse,bp long] (n39) -- (n11);
  \draw[bp stubuse] (n39) -- (n18);
  \draw[bp stubuse] (n39) -- (n21);
  \draw[bp stubuse] (n39) -- (n23);
  \draw[bp stubuse,bp long] (n27) -- (n13);
  \draw[bp stubuse,bp long] (n27) -- (n20);
  \draw[bp stubuse] (n28) -- (n3);
  \draw[bp stubuse] (n28) -- (n5);
  \draw[bp stubuse,bp long] (n28) -- (n21);
  \draw[bp stubuse,bp long] (n28) -- (n23);
  \draw[bp stubuse] (n29) -- (n3);
  \draw[bp stubuse,bp long] (n29) -- (n11);
  \draw[bp stubuse,bp long] (n29) -- (n21);
  \draw[bp stubuse,bp long] (n29) -- (n23);
  \end{scope}
  \begin{scope}[shift={(-0.513,-4.483)}]
    \node[bp R,anchor=west,minimum width=4.5mm] (k1) at (0,0) {R};
    \node[bp R,bp cond,anchor=west,minimum width=4.5mm,right=0.8mm of k1] (k2) {R};
    \node[bp A,anchor=west,minimum width=4.5mm,right=0.8mm of k2] (k3) {A};
    \node[bp P,anchor=west,minimum width=4.5mm,right=0.8mm of k3] (k4) {P};
    \node[bp C,anchor=west,minimum width=4.5mm,right=0.8mm of k4] (k5) {C};
    \node[bp V,anchor=west,minimum width=4.5mm,right=0.8mm of k5] (k6) {V};
    \node[bp key,right=1.2mm of k6] (t1) {严格性标签（虚线框为 $R_{\rm cond}$）};
    \node[bp conj,anchor=west,minimum width=7mm,minimum height=3.4mm,right=2.5mm of t1] (k7) {};
    \node[bp key,right=1mm of k7] (t2) {猜想};
    \node[bp asm,anchor=west,minimum width=4.5mm,right=2.5mm of t2] (k8) {};
    \node[bp key,right=1mm of k8] (t3) {编号假设};
    \node[bp math,anchor=west,minimum width=4.5mm,right=2.5mm of t3] (k9) {};
    \node[bp key,right=1mm of k9] (t4) {纯数学};
    \node[bp stub,anchor=west,minimum width=4.5mm,right=2.5mm of t4] (k10) {};
    \node[bp key,right=1mm of k10] (t5) {其他单元};
    \node[bp fold,anchor=west,minimum width=4.5mm,right=2.5mm of t5] (k11) {$+k$};
    \node[bp key,right=1mm of k11] {折叠};
    \draw[bp use] (0,-0.42) -- ++(6mm,0) coordinate (e1);
    \node[bp key] (u1) at ($(e1)+(1.2mm,0)$) {依赖（前提指向陈述）};
    \draw[bp asmuse] ($(u1.east)+(2.5mm,0)$) -- ++(6mm,0) coordinate (e2);
    \node[bp key] (u2) at ($(e2)+(1.2mm,0)$) {编号假设 $\to R_{\rm cond}$};
    \draw[bp neck] ($(u2.east)+(2.5mm,0)$) -- ++(6mm,0) coordinate (e3);
    \node[bp key] (u3) at ($(e3)+(1.2mm,0)$) {瓶颈依赖（主表）};
    \draw[bp use,bp long] ($(u3.east)+(2.5mm,0)$) -- ++(6mm,0) coordinate (e4);
    \node[bp key] (u4) at ($(e4)+(1.2mm,0)$) {跨两列以上};
    \draw[bp stubuse] ($(u4.east)+(2.5mm,0)$) -- ++(6mm,0) coordinate (e5);
    \node[bp key] at ($(e5)+(1.2mm,0)$) {来自其他单元};
  \end{scope}
\end{tikzpicture}
\caption{第四部分（第~14--18~章）主要定理的依赖图：总表中的陈述与本部分的全部定理、它们在本部分内的直接前提及所用的编号假设（23 个陈述节点、16 个汇总节点）。图例与画法同\cref{fig:appP-dependency}；灰色框代表本部分之外的单元，“$+k$”为按折叠规则省略的 $k$ 个次要陈述。}\label{fig:appP-dependency-e}
\end{figure}
